
\documentclass[11pt,a4paper]{article}

\usepackage[utf8]{inputenc}
\usepackage[T1]{fontenc}
\usepackage{lmodern}
\usepackage{microtype}
\usepackage{amsmath,amssymb,amsthm}
\usepackage{mathtools}
\usepackage{geometry}
\usepackage{hyperref}

\geometry{a4paper,margin=25mm}

\newtheorem{theorem}{Theorem}[section]

\title{\textbf{Identity Bedrock: A Minimal Cryptographic Substrate and Its Systemic Implications}}
\author{Jean-François Bouchard}
\date{March 2026}

\begin{document}

\maketitle

\begin{abstract}
We present Identity Bedrock, a deployed identity system constructed from a minimal cryptographic substrate grounded in one-way functions (OWFs). The system separates identity cryptography from directory consensus via an append-only, chain-agnostic binding layer.

The system has been deployed in production for over five years, with publicly verifiable records and continuous operation.

We demonstrate that naming, content resolution, execution, connectivity, and economic interaction can be derived from this identity substrate without introducing additional trust assumptions.
\end{abstract}

\section{Introduction}

Digital identity systems typically rely on institutional trust anchors such as certificate authorities or identity providers. These introduce external control points and operational dependencies.

This work documents a deployed system that minimizes such dependencies by reducing identity to cryptographic ownership and binding it to an append-only directory.

\section{Background}

One-way functions are necessary and sufficient for constructing secure digital signature schemes. This provides a minimal foundation for authentication.

\section{Architecture}

\subsection{Identity}

Identity is defined as control of a private key $sk$ corresponding to a public key $pk$.

Authentication consists of signing a message $m$:

\begin{equation}
\sigma = \text{Sign}_{sk}(m)
\end{equation}

and verifying:

\begin{equation}
\text{Verify}_{pk}(m, \sigma) = \text{true}
\end{equation}

\subsection{Directory}

An append-only directory binds identifiers to values under a controlling key:

\begin{equation}
(n, v, pk, \sigma)
\end{equation}

where:

\begin{equation}
\sigma = \text{Sign}_{sk}(n \parallel v)
\end{equation}

Verification requires only signature validation.

\section{Security}

\begin{theorem}
If the signature scheme is existentially unforgeable under adaptive chosen-message attack and the directory is append-only with finality, then authentication is unforgeable.
\end{theorem}

\begin{proof}
A successful forgery implies either a break of the signature scheme or a violation of the append-only property of the directory.
\end{proof}

All security properties in the system reduce to the correctness of signature verification and the integrity of the append-only directory.

\section{Derived System Properties}

\subsection{Naming and Resolution}

Naming is implemented through signed bindings:

\begin{equation}
(n, v, pk, \sigma)
\end{equation}

Resolution procedure:

\begin{enumerate}
\item Query name $n$
\item Retrieve $(v, pk, \sigma)$
\item Verify $\sigma$
\item Accept $v$ if verification succeeds
\end{enumerate}

This replaces hierarchical trust with signature verification.

\subsection{Content Addressing}

Content is referenced by cryptographic hash:

\begin{equation}
CID = H(content)
\end{equation}

Integrity holds if recomputed hashes match.

\subsection{Execution}

Execution equivalence is defined by identical inputs, environment specifications, and commands producing identical outputs.

Hash comparison may be used to verify equivalence.

\subsection{Connectivity}

Overlay networking enables communication independent of underlying addressing constraints, without introducing additional trust assumptions.

\subsection{Economic Interaction}

Attribution of actions to identities enables incentive mechanisms without centralized coordination.

\section{Deployment and Operation}

The system is deployed and publicly verifiable.

\begin{center}
\url{https://explorer.emercoin.com/dpo/PrivateNESS.Network}
\end{center}

Verification procedure:

\begin{enumerate}
\item Retrieve record from the directory
\item Extract public key and signature
\item Verify signature
\item Accept binding if valid
\end{enumerate}

Any participant running a compatible node can independently perform this verification.

\section{Implementation}

The system is implemented using Emercoin Name-Value Storage (EmerDNS), anchored via Auxiliary Proof-of-Work (AuxPoW).

Identity bindings are stored as signed records and resolved through DNS-compatible mechanisms (RFC 1035 semantics).

Additional integrations support SSH key distribution and TLS-related bindings without requiring modification to client software.

\section{Evaluation}

Authentication requires only signature verification and directory lookup, resulting in low computational overhead.

The system has operated continuously in production environments for multiple years with publicly verifiable records.

\section{Discussion}

The system demonstrates that naming, resolution, execution, and coordination can be derived from a minimal identity substrate.

No additional trust assumptions are introduced beyond signature security and directory integrity.

\section{Conclusion}

A minimal cryptographic identity substrate, combined with an append-only directory, is sufficient to support a distributed system with verifiable behavior and reduced reliance on centralized infrastructure.

\section*{Acknowledgments}

Portions of this document were developed with the assistance of AI-based language models. All technical design and implementation described herein are the work of the author.

\begin{thebibliography}{9}

\bibitem{rompel1990}
John Rompel.
\textit{One-Way Functions are Necessary and Sufficient for Secure Signatures}.
STOC, 1990.

\bibitem{rfc1035}
P. Mockapetris.
\textit{Domain Names - Implementation and Specification}.
RFC 1035, 1987.

\bibitem{nakamoto2008}
Satoshi Nakamoto.
\textit{Bitcoin: A Peer-to-Peer Electronic Cash System}.
2008.

\bibitem{emercoin}
Emercoin.
\textit{Emercoin NVS and EmerDNS Documentation}.
\url{https://emercoin.com/tech-solutions/}

\end{thebibliography}

\end{document}
